\documentclass[aps,prd,twocolumn,amsmath,amssymb,floatfix]{revtex4-2}

\usepackage{graphicx}
\usepackage{mathtools}
\usepackage{xcolor}

% Macros carried over from the paper
\newcommand{\tAmp}{\widetilde{A}}
\newcommand{\vect}[1]{\ensuremath{\boldsymbol{#1}}}

% Imaginary-part amplitude plots in the thesis layout: A_2B on top, A_12B below,
% each as two coordinate choices. #1 = P_L value (e.g. -1)
\newcommand{\ampstackImThreeD}[1]{%
\begin{tabular}{c}
\includegraphics[width=0.75\textwidth]{../figures/amps/ImA2B_PL#1_3D.pdf}\\
\includegraphics[width=0.75\textwidth]{../figures/amps/ImA2B_Pb_bsq_PL#1_3D.pdf}\\[2ex]
\includegraphics[width=0.75\textwidth]{../figures/amps/ImA12B_PL#1_3D.pdf}\\
\includegraphics[width=0.75\textwidth]{../figures/amps/ImA12B_Pb_bsq_PL#1_3D.pdf}
\end{tabular}}
\newcommand{\ampstackImEta}[1]{%
\begin{tabular}{c}
\includegraphics[width=0.75\textwidth]{../figures/amps/ImA2B_eta8_PL#1_2D.pdf}\\
\includegraphics[width=0.75\textwidth]{../figures/amps/ImA2B_Pb_bsq_eta8_PL#1_2D.pdf}\\[2ex]
\includegraphics[width=0.75\textwidth]{../figures/amps/ImA12B_eta8_PL#1_2D.pdf}\\
\includegraphics[width=0.75\textwidth]{../figures/amps/ImA12B_Pb_bsq_eta8_PL#1_2D.pdf}
\end{tabular}}

% #1 = P_L value, #2 = zeta-hat, #3 = label suffix, #4 = note on the A_12B signal
\newcommand{\imfigset}[4]{%
\begin{figure*}[p]
\centering
\ampstackImThreeD{#1}
\caption{Imaginary parts of the $\tAmp_{2B}$ (top 2 rows) and $\tAmp_{12B}$ (bottom 2 rows) amplitudes for $\vect{P}\cdot a\hat{L}/(2\pi)=\{#1,0,0\}$ ($\hat\zeta=#2$). The first row of each 2-row block displays the $(b_{L}/a,b_{T}/a,\eta|\vect{v}|/a)$ space, while the second row of each 2-row block displays the $((\vect{P}\cdot \vect{b})\cdot\hat{L}/(2\pi),b^{2}/a^2,\eta|\vect{v}|/a)$ space.}
\label{fig:Im_3D_#3}
\end{figure*}
\begin{figure*}[p]
\centering
\ampstackImEta{#1}
\caption{Imaginary parts of the $\tAmp_{2B}$ (top 2 rows) and $\tAmp_{12B}$ (bottom 2 rows) amplitudes for $\vect{P}\cdot a\hat{L}/(2\pi)=\{#1,0,0\}$ ($\hat\zeta=#2$) at fixed $\eta|\vect{v}|/a=8$. The first row of each 2-row block displays the $(b_{L}/a,b_{T}/a)$ space, while the second row of each 2-row block displays the $((\vect{P}\cdot \vect{b})\cdot\hat{L}/(2\pi),b^{2}/a^2)$ space. #4}
\label{fig:Im_eta8_#3}
\end{figure*}}

\begin{document}

\title{Supplemental Material: Dependence of the Sivers TMD Shift on Quark Longitudinal Momentum Fraction $x$ from Lattice QCD}

\author{Hariprashad Ravikumar}
\author{Michael Engelhardt}
\affiliation{Department of Physics, New Mexico State University, Las Cruces, NM 88003, USA}

\maketitle

This Supplemental Material shows the imaginary parts of the extracted amplitudes $\tAmp_{2B}$ and $\tAmp_{12B}$ at the four simulated momenta, $\vect{P}\cdot a\hat{L}/(2\pi)=\{-1,0,0\}$, $\{-2,0,0\}$, $\{-3,0,0\}$ and $\{-4,0,0\}$, corresponding to $\hat\zeta=0.090772$, $0.294367$, $0.539684$ and $0.785756$. For each momentum, the amplitudes are plotted first in the discrete $(b_L/a,b_T/a,\eta|\vect{v}|/a)$ lattice grid and in the Lorentz-invariant $((\vect{P}\cdot\vect{b})\hat{L}/(2\pi),b^2/a^2,\eta|\vect{v}|/a)$ space, and then at fixed $\eta|\vect{v}|/a=8$. The real parts are shown in the main text and in its appendix on the remaining momenta. No signal is discernible in $\tAmp_{12B}^{\text{Im}}$ at any momentum.

\imfigset{-1}{0.090772}{P1}{Notice the lack of discernible signal in $\tAmp_{12B}^{\text{Im}}$.}
\imfigset{-2}{0.294367}{P2}{}
\imfigset{-3}{0.539684}{P3}{}
\imfigset{-4}{0.785756}{P4}{}

\end{document}
